\vspace{-0.8em}
\section{Discussion and Limitations}
\label{sec:disc_lim}
\vspace{-0.5em}

Our method achieves state-of-the-art results across diverse tasks and benchmarks, yet there are several limitations to consider. 
First, our approach relies on leveraging low-level representations from pre-trained, off-the-shelf models. 
Since these representations are frozen during our training, our model's performance is inherently constrained by the quality of these pre-trained embeddings. 
If the pre-trained model fails to effectively capture the content of short video clips, this deficiency may propagate into our long-range understanding framework, affecting overall representation quality. 
However, by choosing high-performing video-text alignment models such as LanguageBind, we mitigate this limitation and ensure strong baseline representations.

Second, unlike methods operating on image patches, where masked token predictions can be easily visualized, our approach operates on embeddings, which are inherently more abstract and less interpretable.
This makes it challenging to directly assess the model's performance on masked token predictions. 
To address this, we propose an interpretability strategy that leverages captions or generates them when unavailable, providing a proxy for visualizing and understanding the model’s predictions. 
This approach aids in interpreting predictions, though it remains an approximation rather than a direct visualization of learned embeddings.